% Einheit 2 von 2 – Spiegelung an der Winkelhalbierenden (bank/umkehrfunktion/e2.jsonl, zusammenbau v0.1)
%% TODO Zeitmarke Einheit 2: Marken-Zeile nicht lesbar
%% TODO Zweigzeile Teil 1: was hier gelernt wird (Satz in Schülersprache aus dem Einheitstitel) – Einheit 2
\einheitenkopf[e2]{Einheit 2 von 2 · Spiegelung an der Winkelhalbierenden}
\zweigzeile{Abitur LK}
\setcounter{aufgabe}{10}

%% TODO Titel: Ich-kann-Satz fehlt in der Bank (Platzhalter: Kettenname) – Nr. 11
%% TODO Anweisung: ein Satz über der ersten Teilaufgabe fehlt in der Bank – Nr. 11
\begin{aufgabe}{Spiegelung an der Winkelhalbierenden}
%% TODO \rechenplatz{n}: Schrittzahl des Grundfalls fehlt in der Bank (nur bei fester Schreibform, 2.3 b)
\begin{teile}
\teil Der Punkt (2 | 5) wird an der Geraden $y = x$ gespiegelt. Wo liegt das Bild? \kreuz{$(5 | 2)$} \\ \kreuz{$(-2 | 5)$} \\ \kreuz{$(2 | -5)$}
\teil Spiegle die Punkte A(2 | 6) und B(−1 | 4) an der Geraden $y = x$.
\teil Die Tangente an den Graphen von f im Punkt P(1 | 2) hat die Steigung 3. g ist die Umkehrfunktion. Welche Steigung hat die Tangente an den Graphen von g im Punkt P'(2 | 1)? \leerfeld
\teil $f(x) = 6e^{-\frac{1}{2}x} - 2$ ist umkehrbar, g ist die Umkehrfunktion. t ist die Tangente an den Graphen von f im Punkt Q mit der Steigung −1. Begründe ohne Rechnung, dass t auch den Graphen von g berührt, und gib den Berührpunkt an \hfill (Abitur 2026 LK).
\teil Die Punkte A(0 | 4) und B(1 | 7) werden an $y = x$ gespiegelt. Begründe, dass A, B, B' und A' ein Trapez bilden, und berechne seinen Flächeninhalt. \leerfeld
\teil f und g sind auf $[0;\, \infty[$ definiert, g ist die Umkehrfunktion von f. Die Abbildung zeigt beide Graphen; sie schneiden sich im Ursprung und im Punkt S mit der x-Koordinate $x_S$. Beurteile ohne Rechnung die Aussage: „Das Integral von 0 bis $x_S$ über $g(x) - x$ ist gleich dem Integral von 0 bis $x_S$ über $x - f(x)$.“ \hfill (Abitur 2025 LK)

\begin{ksys}[xmin=0,xmax=5,ymin=0,ymax=5,ablesen] \funktionab{0.25*\x^2}{f}{0}{4.4} \funktionab{2*sqrt(\x)}{g}{0}{5} \gerade{1}{0}{y = x} \end{ksys}
\end{teile}
\end{aufgabe}

%% TODO Titel: Ich-kann-Satz fehlt in der Bank (Platzhalter: Kettenname) – Nr. 12
%% TODO Anweisung: ein Satz über der ersten Teilaufgabe fehlt in der Bank – Nr. 12
\begin{aufgabe}{Spiegelung an der Winkelhalbierenden – Fehler finden, Begründen, Anwendung, Darstellungswechsel}
\begin{teile}
\teil Mia sagt: „Die Fläche zwischen dem Graphen von f und der Geraden y = x ist sicher kleiner als die zwischen y = x und dem Graphen der Umkehrfunktion, weil die Graphen verschieden aussehen.“ Finde den Fehler.
\teil Begründe, warum die Spiegelung an der Geraden y = x die Koordinaten eines Punktes vertauscht.
\teil Eine Pflanze wächst am fünften Tag um 2 cm pro Tag. Wie viele Tage braucht sie an dieser Stelle je Zentimeter Wachstum? Deute die Antwort als Steigung der Umkehrfunktion.
\teil Die Abbildung zeigt den Graphen der Umkehrfunktion g. Lies drei Punkte ab und gib die zugehörigen Punkte des Graphen von f an.

\begin{ksys}[xmin=-1,xmax=6,ymin=-3,ymax=6,ablesen] \funktionab{ln(\x)/ln(2)}{g}{0.1}{5.5} \gerade{1}{0}{y = x} \end{ksys}
\end{teile}
\end{aufgabe}
